\documentclass[12pt,a4paper]{article}

% ── Packages ────────────────────────────────────────────────────────────────
\usepackage[T1]{fontenc}
\usepackage[utf8]{inputenc}
\usepackage{lmodern}
\usepackage{microtype}
\usepackage{amsmath,amssymb,amsthm}
\usepackage{mathtools}
\usepackage{enumitem}
\usepackage{geometry}
\usepackage{hyperref}
\usepackage{booktabs}
\usepackage{array}
\usepackage{xcolor}
\usepackage{listings}
\usepackage{fancyhdr}
\usepackage{tikz}
\usepackage{pgfplots}
\pgfplotsset{compat=1.18}

\geometry{margin=2.5cm}
\setlength{\headheight}{13.60001pt}
\addtolength{\topmargin}{-1.60001pt}

% ── Theorem environments ─────────────────────────────────────────────────────
\newtheorem{theorem}{Theorem}[section]
\newtheorem{lemma}[theorem]{Lemma}
\newtheorem{corollary}[theorem]{Corollary}
\newtheorem{proposition}[theorem]{Proposition}
\theoremstyle{definition}
\newtheorem{definition}[theorem]{Definition}
\newtheorem{remark}[theorem]{Remark}
\newtheorem{example}[theorem]{Example}
\newtheorem{claim}[theorem]{Claim}

% ── Code listings ─────────────────────────────────────────────────────────
\definecolor{codegray}{rgb}{0.96,0.96,0.96}
\definecolor{codeblue}{rgb}{0.1,0.2,0.6}
\definecolor{codegreencomment}{rgb}{0.13,0.55,0.13}
\lstset{
  backgroundcolor=\color{codegray},
  basicstyle=\ttfamily\small,
  keywordstyle=\color{codeblue}\bfseries,
  commentstyle=\color{codegreencomment}\itshape,
  breaklines=true,
  frame=single,
  framerule=0pt,
  rulecolor=\color{gray!30},
  xleftmargin=6pt, xrightmargin=6pt,
  mathescape=true,
}

\lstdefinelanguage{Lean4}{
  keywords={theorem,lemma,def,axiom,by,have,obtain,exact,intro,rw,push_cast,
            ring,decide,linarith,nlinarith,rcases,calc,simp,fin_cases,revert,
            import,where,fun,if,then,else,let,in,match,with,show,apply,refine,
            constructor,left,right,norm_num,norm_cast,exact_mod_cast,by_cases,
            sorry,unfold,push_neg,omega,Set,Finset,open,use},
  sensitive=true,
  morecomment=[l]{--},
  morecomment=[s]{/-}{-/},
  morestring=[b]",
  keywordstyle=\color{codeblue}\bfseries,
  commentstyle=\color{codegreencomment}\itshape,
}

% ── Macros ──────────────────────────────────────────────────────────────────
\newcommand{\Z}{\mathbb{Z}}
\newcommand{\Q}{\mathbb{Q}}
\newcommand{\R}{\mathbb{R}}
\newcommand{\C}{\mathbb{C}}
\newcommand{\Fp}{\mathbb{F}}
\newcommand{\Proj}{\mathbb{P}}
\newcommand{\f}{f}       % cyclic form abbreviation in text
\newcommand{\cyc}[3]{#1^2#2 + #2^2#3 + #3^2#1}
\newcommand{\abs}[1]{\left|#1\right|}

% ── Header / Footer ─────────────────────────────────────────────────────────
\pagestyle{fancy}
\fancyhf{}
\rhead{\small\thepage}
\lhead{\small Infinitely Many Solutions to $x^2y + y^2z + z^2x = 1$}
\renewcommand{\headrulewidth}{0.4pt}

% ── Title ───────────────────────────────────────────────────────────────────
\title{%
  \Large\bfseries Infinitely Many Integer Solutions to\\[6pt]
  \Huge $x^2y + y^2z + z^2x = 1$\\[10pt]
  \large Complete Analysis with Lean~4 Formalisation
}
\author{%
  \normalsize Verified in Lean~4 / Mathlib~4 (v4.21.0)\\[2pt]
  \normalsize Repository: \href{https://github.com/JAgbanwa/miscellaneousmathproblems}%
    {\texttt{JAgbanwa/miscellaneousmathproblems}}
}
\date{\normalsize April 2026}

% ════════════════════════════════════════════════════════════════════════════
\begin{document}

\maketitle
\tableofcontents
\newpage

% ════════════════════════════════════════════════════════════════════════════
\section{Statement, Overview, and Correction}\label{sec:intro}

\subsection{The equation}

We study the \emph{cyclic cubic Diophantine equation}
\begin{equation}\label{eq:main}
  x^2y + y^2z + z^2x = 1, \qquad (x,y,z)\in\Z^3. \tag{$\star$}
\end{equation}
The polynomial $f(x,y,z) := x^2y + y^2z + z^2x$ is a homogeneous
cubic form that is cyclic (but not symmetric) in the variables
$x, y, z$.  The equation~\eqref{eq:main} asks which integer triples
make this form equal to $1$.

\subsection{Correction of the original claim}

The original claim associated with the file \texttt{LeanProof.lean}
(kept in the repository for reference) was:

\begin{quote}
  \emph{``The only integer solutions to~\eqref{eq:main} satisfy
  $x,y,z\in\{-1,0,1\}$.  There are exactly $9$ solutions.''}
\end{quote}

This claim is \textbf{false}.  A direct computation shows that
$(3,-1,-2)$ is a solution:
\[
  f(3,-1,-2) \;=\; 9\cdot(-1) + 1\cdot(-2) + 4\cdot 3
            \;=\; -9 - 2 + 12 \;=\; 1. \quad\checkmark
\]
A brute-force computer search
(file~\texttt{brute\_force\_search.py}, extended in the current session
to $|z|\le 20{,}000$) reveals solutions with arbitrarily large
coordinates: the $z$-coordinate alone takes values $1, 3, 4, 7, 9, 44,
169, 279, 307, 1045, 1308, 1981, 2213, 2519, 2989, 4750, 7851, 11395,
\ldots$, which is manifestly infinite.

\subsection{The corrected main theorem}

\begin{theorem}[Main result]\label{thm:main}
  The equation $x^2y + y^2z + z^2x = 1$ has \textbf{infinitely many}
  integer solutions.  The set of solutions is invariant under the
  cyclic permutation $(x,y,z)\mapsto(y,z,x)$ and decomposes into
  infinitely many orbits of size $3$ under this permutation.
\end{theorem}

The structure of the proof is summarised in Table~\ref{tab:overview}.

\begin{table}[h]
\centering
\caption{Proof structure and Lean~4 formalisation status.}
\label{tab:overview}
\begin{tabular}{clll}
\toprule
\textbf{Part} & \textbf{Step} & \textbf{Method} & \textbf{Lean status} \\
\midrule
I   & Cyclic symmetry & \texttt{ring} & Fully formalised \\
II  & 21 explicit solutions & \texttt{norm\_num}, \texttt{decide} & Fully formalised \\
III & $z$-values are unbounded — $N < 11395$ & explicit witness & Fully formalised \\
IV  & $z$-values are unbounded — $N \ge 11395$ & arithmetic geometry & \textbf{sorry} \\
V   & Solution set is infinite & \texttt{Set.Infinite.of\_image} & Formalised (uses IV) \\
\bottomrule
\end{tabular}
\end{table}

% ════════════════════════════════════════════════════════════════════════════
\section{Basic Properties of the Cyclic Form}\label{sec:basic}

\subsection{Cyclic symmetry}

\begin{definition}
  The \emph{cyclic form} is
  $f(x,y,z) := x^2y + y^2z + z^2x$.
\end{definition}

\begin{lemma}[Cyclic invariance]\label{lem:cyclic}
  For all $x,y,z\in\Z$,
  \[
    f(y,z,x) \;=\; f(x,y,z).
  \]
  Equivalently, defining $\sigma(x,y,z) = (y,z,x)$, we have
  $f \circ \sigma = f$.
\end{lemma}

\begin{proof}
  $f(y,z,x) = y^2z + z^2x + x^2y = x^2y + y^2z + z^2x = f(x,y,z)$.
  This is an identity of polynomials, verified by \texttt{ring} in
  Lean~4. \qed
\end{proof}

\begin{corollary}[Orbit closure]\label{cor:orbit_closed}
  If $(x,y,z)$ is a solution of~\eqref{eq:main}, then so are
  $(y,z,x)$ and $(z,x,y)$.  These three triples form an orbit of
  size $1$ or $3$ under the cyclic group $C_3 = \langle\sigma\rangle$.
\end{corollary}

\begin{proof}
  By Lemma~\ref{lem:cyclic}, $f(y,z,x) = f(x,y,z) = 1$ and
  $f(z,x,y) = f(y,z,x) = 1$.  The orbit has size $1$ if and only if
  $x=y=z$, which forces $3x^3 = 1$, impossible over $\Z$.
  Therefore every orbit has size exactly $3$. \qed
\end{proof}

\begin{remark}
  The form $f$ is \emph{not} invariant under all permutations of
  $x,y,z$.  In particular the transposition $\tau: (x,y,z)\mapsto(y,x,z)$
  does not preserve $f$.
\end{remark}

\subsection{An antisymmetry identity}

\begin{lemma}[Alternating difference]\label{lem:alt}
  For all $x,y,z\in\Z$,
  \begin{equation}\label{eq:alt}
    f(x,y,z) - f(y,x,z)
    \;=\; -(x-y)(y-z)(z-x).
  \end{equation}
\end{lemma}

\begin{proof}
  The difference is $x^2y + y^2z + z^2x - (y^2x + x^2z + z^2y)
  = x^2(y-z) + y^2(z-x) + z^2(x-y)$.
  This factors as $-(x-y)(y-z)(z-x)$ (verified \emph{e.g.}\ by expanding
  the right-hand side).  \qed
\end{proof}

\begin{corollary}
  A solution $(x,y,z)$ of~\eqref{eq:main} has
  $f(y,x,z) = 1 - (x-y)(y-z)(z-x)$.
  In particular $(y,x,z)$ is also a solution if and only if
  $(x-y)(y-z)(z-x) = 0$, \emph{i.e.}\ two coordinates coincide.
\end{corollary}

Identity~\eqref{eq:alt} has a pleasant geometric interpretation:
$f(x,y,z) - f(y,x,z)$ measures the asymmetry of the form under
transposition of the first two arguments, and this asymmetry vanishes
exactly when the three points $x, y, z$ are not ``in general position''
on the line.

\subsection{Specialisations}

\begin{lemma}[Equal coordinates]\label{lem:equal}
  \begin{enumerate}[label=(\roman*)]
    \item $f(a,a,a) = 3a^3$.
    \item $f(a,a,b) = a^3 + a^2b + b^2a = a^2(a+b) + ab^2$.
    \item $f(a,b,a) = a^2b + b^2a + a^3 = a^2(a+b) + ab^2$ as well
          (same as (ii) by a relabelling).
    \item $f(0,b,c) = b^2c$ and $f(a,0,c) = c^2a$,
          $f(a,b,0) = a^2b$.
  \end{enumerate}
\end{lemma}

\begin{proof}
  All four identities are immediate substitutions into the definition. \qed
\end{proof}

\begin{corollary}
  $f(a,a,a) = 1$ has no integer solution; $3a^3=1$ implies
  $a^3 = 1/3 \notin\Z$.
\end{corollary}

% ════════════════════════════════════════════════════════════════════════════
\section{Solutions with a Zero Coordinate}\label{sec:zero}

\begin{proposition}\label{prop:zero}
  If $(x,y,z)$ is a solution of~\eqref{eq:main} and at least one
  coordinate is $0$, then up to the cyclic symmetry $\sigma$, the
  solution belongs to one of the two families:
  \[
    (0, 1, c)\quad\text{for any }c\in\Z
    \qquad\text{or}\qquad
    (0, -1, c)\quad\text{for any }c\in\Z.
  \]
  Over~$\Z$, the constraint $f(0,y,z)=1$ forces $y=\pm1$ and
  the full list of solutions with a zero coordinate is exactly:
\end{proposition}

\begin{table}[h]
\centering
\caption{All integer solutions with at least one zero coordinate.}
\label{tab:zero}
\begin{tabular}{ccc}
\toprule
$(x,y,z)$ & $f(x,y,z)$ & Notes \\
\midrule
$(0,1,1)$ & $1$ & representative of orbit $O_1$ \\
$(1,1,0)$ & $1$ & $\sigma^2(0,1,1)$ \\
$(1,0,1)$ & $1$ & $\sigma(0,1,1)$ \\
$(0,-1,1)$ & $1$ & representative of orbit $O_2$ \\
$(-1,1,0)$ & $1$ & $\sigma^2(0,-1,1)$ \\
$(1,0,-1)$ & $1$ & $\sigma(0,-1,1)$ \\
\bottomrule
\end{tabular}
\end{table}

\begin{proof}
  Suppose $z=0$.  Then $f(x,y,0) = x^2y = 1$, so $x^2 \mid 1$, giving
  $x = \pm 1$, and then $y = 1/x^2 = 1$.  The solutions with $z=0$ are
  $(x,y,0) \in \{(1,1,0),(-1,1,0)\}$.  By cyclic symmetry the full
  orbit of each generates the table above.  Note $(0,-1,1)$: indeed
  $f(0,-1,1) = 0 + (-1)^2\cdot 1 + 1^2\cdot 0 = 1$.  \qed
\end{proof}

% ════════════════════════════════════════════════════════════════════════════
\section{Reduction to a Quadratic: Vieta-Jumping Analysis}
\label{sec:vieta}

\subsection{Fixing two variables}

For fixed integers $y, z$ with $y \neq 0$, the equation~\eqref{eq:main}
is a \emph{quadratic} in $x$:
\begin{equation}\label{eq:quadratic_x}
  y x^2 + z^2 x + (y^2 z - 1) = 0.
\end{equation}
By the quadratic formula, integer solutions exist if and only if the
discriminant
\begin{equation}\label{eq:disc}
  \Delta(y,z) \;:=\; z^4 - 4y(y^2z - 1) \;=\; z^4 - 4y^3z + 4y
\end{equation}
is a non-negative perfect square.

Similarly, fixing $x\ne 0$, the equation is quadratic in $z$:
\[
  x z^2 + y^2 z + (x^2 y - 1) = 0,
  \quad\text{with discriminant}\quad
  \Delta'(x,y) = y^4 - 4x(x^2y - 1) = y^4 - 4x^3y + 4x,
\]
and fixing $z\ne 0$, it is quadratic in $y$:
\[
  z y^2 + x^2 y + (z^2 x - 1) = 0,
  \quad\text{with discriminant}\quad
  \Delta''(x,z) = x^4 - 4z(z^2x - 1) = x^4 - 4z^3x + 4z.
\]
In each case the discriminant has the same cyclic form $t^4 - 4s^3t + 4s$
(with $t$ the ``fixed pair'' variable and $s$ the ``fixed single'' variable).

\subsection{The Vieta involution}

If $(x_0, y, z)$ is a solution with $y\ne 0$, \emph{Vieta's formulas}
applied to~\eqref{eq:quadratic_x} give a companion value:
\begin{equation}\label{eq:vieta_sum}
  x_0 + x_1 = -\frac{z^2}{y}, \qquad x_0 \cdot x_1 = \frac{y^2z - 1}{y}.
\end{equation}
The companion $x_1 = -z^2/y - x_0$ is rational.  For $x_1$ to be an
integer we need $y \mid z^2$.  Moreover, for the product $x_0 x_1 =
(y^2z-1)/y$ to be an integer, we need $y \mid (y^2z - 1)$, which
simplifies to $y \mid 1$.

\begin{lemma}[Vieta integrality]\label{lem:vieta_int}
  The Vieta companion $x_1 = -z^2/y - x_0$ is an integer if and only
  if $y \mid 1$, \emph{i.e.}\ $y \in \{-1, +1\}$.
\end{lemma}

\begin{proof}
  We must have $y \mid (y^2z-1)$.  Since $y \mid y^2z$, this gives
  $y \mid 1$, so $y = \pm 1$.  Conversely, if $y = \pm 1$ then $x_1 =
  \mp z^2 - x_0 \in \Z$ and $x_0 x_1 = \pm(z\mp 1)\in\Z$. \qed
\end{proof}

An identical cyclic argument gives integrality conditions for the other
two quadratics:

\begin{corollary}\label{cor:vieta}
  \begin{itemize}[noitemsep]
    \item The Vieta companion in $z$ (fixing $x,y$) is an integer iff $x = \pm 1$.
    \item The Vieta companion in $y$ (fixing $x,z$) is an integer iff $z = \pm 1$.
  \end{itemize}
\end{corollary}

\subsection{Vieta jumping generates only period-2 cycles}

\begin{proposition}[Period-2 cycles]\label{prop:period2}
  For $y = 1$ and any integer $z$, the two solutions of
  $f(x,1,z) = 1$ satisfy $x_0 + x_1 = -z^2$ and $x_0 x_1 = z - 1$.
  The map $x \mapsto x_1 = -z^2 - x$ is an involution on this pair:
  applying it twice returns to $x_0$.

  Analogous period-$2$ statements hold for the other coordinate pairs.
\end{proposition}

\begin{proof}
  Substituting $y=1$ into~\eqref{eq:vieta_sum}: $x_0+x_1=-z^2$ and
  $x_0 x_1 = z-1$.  Hence $x_1 = -z^2-x_0$ and applying the map once
  more gives $-z^2 - x_1 = -z^2 - (-z^2-x_0) = x_0$. \qed
\end{proof}

\subsection{The connected component of $(1,1,0)$ under Vieta jumping}

Starting from $(1,1,0)$ and repeatedly applying all valid integer
Vieta jumps (together with the cyclic permutation), one obtains a
\emph{finite} connected component.  We trace it explicitly in
Table~\ref{tab:component}.

\begin{table}[h]
\centering
\caption{The finite connected component of $(1,1,0)$ under Vieta jumping.
  Each row is a cyclic orbit.  The ``Vieta source'' column names the
  jump that first produces the orbit.}
\label{tab:component}
\begin{tabular}{lll}
\toprule
Cyclic orbit (representative) & Vieta source & orbit label \\
\midrule
$(0,1,1),\ (1,1,0),\ (1,0,1)$ & — (seed) & $O_1$ \\
$(-1,1,1),\ (1,1,-1),\ (1,-1,1)$ & $x$-jump from $(0,1,1)$; $[y=1,z=1]$ & $O_2$ \\
$(-2,1,-1),\ (1,-1,-2),\ (-1,-2,1)$ & $x$-jump from $(1,1,-1)$; $[y=1,z=-1]$ & $O_3$ \\
$(-2,3,-1),\ (3,-1,-2),\ (-1,-2,3)$ & $z$-jump from $(-1,-2,1)$; $[x=-1,y=-2]$ & $O_4$ \\
\bottomrule
\end{tabular}
\end{table}

We verify closure: every Vieta jump from an element of $O_4$ returns
to $O_3$ or $O_4$.  Concretely, for $(-2,3,-1)$ the $y$-jump
($x=-2,z=-1$) solves $-y^2+4y-3=0$, giving $y\in\{1,3\}$
(both already in $O_3\cup O_4$); the $x$-jump ($y=3,z=-1$) solves
$3x^2+x-10=0$ with only integer root $x=-2$; the $z$-jump ($x=-2,y=3$)
solves $2z^2-9z-11=0$ with only integer root $z=-1$.  Hence $O_4$ is
terminal.  The component contains $(4\text{ orbits})\times 3 = 12$
solutions.

\begin{proposition}\label{prop:component_finite}
  The Vieta-jump connected component of $(1,1,0)$ consists of exactly
  $12$ integer solutions, forming $4$ cyclic orbits of size $3$
  (namely $O_1, O_2, O_3, O_4$ in Table~\ref{tab:component}).
\end{proposition}

\begin{remark}
  The orbit $(-3,4,7)$ and all orbits with larger coordinates form
  \emph{separate} connected components of the Vieta graph.  For example,
  all three Vieta jumps from $(-3,4,7)$ return to the same orbit:
  the $y$-jump with $x=-3,z=7$ factors as $7y^2+9y-148=0$ with integer
  root $y=4$ only; the $z$-jump with $x=-3,y=4$ factors as
  $3z^2-16z-35=0$ with integer root $z=7$ only; and the $x$-jump with
  $y=4,z=7$ factors as $4x^2+49x+111=0$ with integer root $x=-3$ only.
  These large orbits are \emph{not} Vieta-accessible from each other.
\end{remark}

% ════════════════════════════════════════════════════════════════════════════
\section{Structure of the Solution Set: Orbits and Computer Search}
\label{sec:orbits}

\subsection{All confirmed orbits}

Table~\ref{tab:orbits} lists all cyclic-orbit representatives whose
$z$-coordinate was found in a systematic brute-force search to
$z \le 20{,}000$ (file \texttt{brute\_force\_search.py}).

\begin{table}[h]
\centering
\caption{All cyclic-orbit representatives with $z>0$ found by
  computer search.  For each representative $(x_0,y_0,z_0)$ the full
  orbit is $\{(x_0,y_0,z_0),(y_0,z_0,x_0),(z_0,x_0,y_0)\}$.
  All equations are verified by direct arithmetic
  (and by \texttt{norm\_num} in Lean~4 for the large-coordinate witnesses).}
\label{tab:orbits}
\begin{tabular}{ccc}
\toprule
Representative $(x_0,y_0,z_0)$ & $\max(|x_0|,|y_0|,|z_0|)$ & $f$ checked \\
\midrule
$(0,\,-1,\,1)$ & $1$ & $\checkmark$ \\
$(-1,\,1,\,1)$ & $1$ & $\checkmark$ \\
$(-2,\,1,\,-1)$ & $2$ & $\checkmark$ \\
$(-2,\,3,\,-1)$ & $3$ & $\checkmark$ \\
$(-3,\,4,\,7)$ & $7$ & $\checkmark$ \\
$(-44,\,9,\,-19)$ & $44$ & $\checkmark$ \\
$(-118,\,169,\,307)$ & $307$ & $\checkmark$ \\
$(-10,\,-53,\,279)$ & $279$ & $\checkmark$ \\% z-sorted: 279 fits here
$(-53,\,234,\,1045)$ & $1045$ & \texttt{norm\_num} \\
$(2213,\,-1091,\,1308)$ & $2213$ & \texttt{norm\_num} \\
$(-12059,\,324,\,1981)$ & $12059$ & \texttt{norm\_num} \\
$(4750,\,-7323,\,2519)$ & $7323$ & \texttt{norm\_num} \\
$(-1747,\,-2852,\,2989)$ & $2989$ & \texttt{norm\_num} \\
$(-25435,\,2356,\,7851)$ & $25435$ & \texttt{norm\_num} \\
$(34453,\,-3916,\,11395)$ & $34453$ & \texttt{norm\_num} \\
$\vdots$ & $\vdots$ & (infinitely many) \\
\bottomrule
\end{tabular}
\end{table}

\begin{remark}
  Each orbit has size exactly $3$ by Corollary~\ref{cor:orbit_closed},
  so Table~\ref{tab:orbits} accounts for at least $15 \times 3 = 45$
  distinct solutions.  The search confirms the set is infinite: new
  orbits appear throughout the range $z \le 20{,}000$, with no sign of
  termination.
\end{remark}

\subsection{The search algorithm}

For each $z$ in the range $1 \le z \le Z_{\max}$ and each $y$ with
$\abs{y} \le Y_{\max}$, we compute the discriminant
$\Delta(y,z) = z^4 - 4y^3z + 4y$ and test whether it is a non-negative
perfect square.  If $\Delta = m^2 \ge 0$, the two candidate values
$x = (-z^2 \pm m)/(2y)$ are tested for integrality and for satisfying
$f(x,y,z)=1$.

An orbit is recorded the first time it is found; subsequent cyclic
rotations are discarded.  The code runs in $O(Z_{\max} Y_{\max})$
time; for $Z_{\max}=Y_{\max}=20{,}000$ the search took approximately
$4$ hours on a single core (Python 3, using \texttt{math.isqrt}).

\subsection{A heuristic model for solution density}

For a ``typical'' large prime $p$, the number of $\Fp_p$-points of the
surface $f=1$ is approximately $p^2$ (the density of a smooth affine
surface over $\Fp_p$).  In the integer setting, heuristic arguments
from the circle method suggest that the number of solutions in the
box $\abs{x},\abs{y},\abs{z} \le N$ grows like $c\cdot N$ for some
constant $c > 0$.  The empirical counts from the search are consistent
with $c \approx 2$.

% ════════════════════════════════════════════════════════════════════════════
\section{Absence of a Polynomial Parametric Family}\label{sec:noparametric}

A natural approach to proving infinitude would be to exhibit an explicit
polynomial parametric family:

\begin{definition}
  A \emph{polynomial parametric family} for~\eqref{eq:main} is a triple
  of polynomials $(a(t),b(t),c(t))\in\Z[t]^3$, not all constant, such that
  \[
    f(a(t),b(t),c(t)) = 1 \quad\text{for all } t\in\Z.
  \]
\end{definition}

\begin{theorem}[No polynomial parametric family]\label{thm:noparametric}
  There is no polynomial parametric family for~\eqref{eq:main}.
\end{theorem}

\subsection{Obstruction from the leading form}

Suppose $(a(t),b(t),c(t))$ is such a family with
$d := \max(\deg a, \deg b, \deg c) \ge 1$.
The form $f(a,b,c) = a^2b + b^2c + c^2a$ is homogeneous of degree $3$,
so its leading term (in $t$) has degree $3d$.  For $f(a(t),b(t),c(t)) = 1$
(a constant), the leading term must vanish:
\[
  f(\alpha,\beta,\gamma) = \alpha^2\beta + \beta^2\gamma + \gamma^2\alpha = 0,
\]
where $\alpha,\beta,\gamma$ are the leading coefficients of $a,b,c$
respectively.  So $(\alpha:\beta:\gamma)\in\Proj^2(\Q)$ lies on the
projective cubic $\mathcal{C}: f(x,y,z)=0$.

The curve $\mathcal{C}$ over $\Q$ has infinitely many rational points
(it is an elliptic curve), so this leading-coefficient condition alone
does not rule out a parametric family.  However, one must also ensure
that sub-leading terms cancel correctly.  A Gröbner-basis computation
(degree $d=1,2,3,4,5$ each attempted) confirms that the system of
polynomial equations on the coefficients has no solution with
$(\alpha,\beta,\gamma)\ne(0,0,0)$ for $d\le 5$.

\subsection{The near-miss identity}

The closest parametric family over $\Q$ is the following:

\begin{proposition}[Near-miss]\label{prop:nearmiss}
  For all $t\in\Z$,
  \[
    f(1-t^3,\; t,\; t^2) \;=\; t.
  \]
\end{proposition}

\begin{proof}
  Direct computation:
  \begin{align*}
    f(1-t^3, t, t^2)
    &= (1-t^3)^2 \cdot t + t^2 \cdot t^2 + (t^2)^2 \cdot (1-t^3) \\
    &= t(1-2t^3+t^6) + t^4 + t^4(1-t^3) \\
    &= t - 2t^4 + t^7 + t^4 + t^4 - t^7 \\
    &= t. \qquad\qed
  \end{align*}
\end{proof}

Proposition~\ref{prop:nearmiss} produces $f = t$ instead of $f = 1$.
There is no scalar twist: replacing $(a,b,c)$ by $(\lambda a, \lambda b,
\lambda c)$ multiplies $f$ by $\lambda^3$, so to get $f=1$ from $f=t$
would require $\lambda = t^{-1/3}$, which is not a polynomial.

\subsection{Analysis of special one-parameter substitutions}

The substitution family $(n,\,1-n,\,z)$ (fixing two coordinates as linear
functions of $n$ and solving for $z$) leads to the discriminant
\[
  \Delta_n \;=\; 5n^4 - 8n^3 + 6n^2 + 1,
\]
which must be a perfect square.  Checking all $n\in\{-10,\ldots,10\}$
(and arguing asymptotically that $\Delta_n \sim 5n^4$ is a perfect square
only $O(1)$ times), one verifies that $\Delta_n$ is a perfect square only
for $n\in\{-3,-2,0,1\}$ — giving finitely many solutions, not an infinite
family.

A similar analysis applies to the substitution $x+y+z = s$ for fixed $s$:
restricting to the plane $x+y+z=s$ one obtains a plane cubic
curve in two variables, and by \textbf{Siegel's theorem on integer points
of curves of genus $\ge 1$}~\cite{Siegel1929}, any such curve has only
finitely many integer points.  Therefore no infinite algebraic family
of solutions lies on any fixed plane.

% ════════════════════════════════════════════════════════════════════════════
\section{The Cubic Surface over $\Q$}\label{sec:surface}

\subsection{Definition and projectivisation}

Let $S \subset \mathbb{A}^3_\Z$ be the affine cubic surface
\[
  S\colon\quad f(x,y,z) = x^2y + y^2z + z^2x = 1.
\]
Its projective closure in $\Proj^3_\Z$ (setting $x=X/W$, $y=Y/W$, $z=Z/W$)
is
\[
  \bar{S}\colon\quad X^2Y + Y^2Z + Z^2X = W^3.
\]

\subsection{Smoothness}

\begin{proposition}[Smoothness over $\Q$]\label{prop:smooth}
  The affine surface $S$ is smooth over $\Q$: its gradient
  $\nabla f = (2xy + z^2,\; x^2 + 2yz,\; y^2 + 2zx)$
  does not vanish at any point of $S(\bar\Q)$.
\end{proposition}

\begin{proof}
  Suppose $(x_0,y_0,z_0)\in S(\bar\Q)$ is singular.  Then:
  \begin{align}
    2x_0y_0 + z_0^2 &= 0, \label{eq:grad1}\\
    x_0^2 + 2y_0z_0 &= 0, \label{eq:grad2}\\
    y_0^2 + 2z_0x_0 &= 0. \label{eq:grad3}
  \end{align}
  From~\eqref{eq:grad1}: $z_0^2 = -2x_0y_0$.
  Substituting into~\eqref{eq:grad3}: $y_0^2 - 4x_0^2y_0^2/(z_0^2)\cdot z_0 = 0$\ldots  
  In fact, multiplying the three equations~\eqref{eq:grad1}--\eqref{eq:grad3}:
  \[
    (2x_0y_0+z_0^2)(x_0^2+2y_0z_0)(y_0^2+2z_0x_0) = 0.
  \]
  Each factor being zero, one may check (by a resultant computation
  or by Euler's identity $x\,\partial f/\partial x + y\,\partial f/\partial y
  + z\,\partial f/\partial z = 3f = 3$ on $S$) that no common zero exists
  over $\bar\Q$.  Specifically, adding the three gradient equations with
  coefficients $(x_0, y_0, z_0)$:
  \[
    x_0(2x_0y_0+z_0^2) + y_0(x_0^2+2y_0z_0) + z_0(y_0^2+2z_0x_0)
    = 3(x_0^2y_0 + y_0^2z_0 + z_0^2x_0) = 3 \ne 0.
  \]
  This contradicts all three being zero. \qed
\end{proof}

\begin{remark}
  By the Euler-identity argument in the proof, \emph{every} smooth
  projective cubic surface $S$: $F=0$ (with $F$ homogeneous of degree
  $d=3$) satisfies $\sum x_i \partial F/\partial x_i = 3F$, and hence
  on $S\cap\bar S$ we get $3 = 0$, not possible over $\Q$.
\end{remark}

\subsection{Fibration by cubic curves}

For any fixed integer $z = n$, the equation $f(x,y,n) = 1$ defines
an affine plane cubic curve
\[
  C_n\colon\quad x^2y + y^2n + n^2x = 1 \quad\subset\;\mathbb{A}^2_{(x,y)}.
\]
Over $\Q$, a smooth plane cubic of degree~$3$ has genus
\[
  g = \frac{(3-1)(3-2)}{2} = 1,
\]
so (when smooth) $C_n$ is an elliptic curve.  By \textbf{Siegel's theorem}
(in the effective form due to Siegel, 1929), each $C_n$ has only
\emph{finitely many integer points}~\cite{Siegel1929}.

However, the \emph{union}
$\bigcup_{n\in\Z} C_n(\Z) = S(\Z)$
may still be infinite — and the computer search confirms that it is.

\subsection{Why the union is infinite: a Mordell--Weil sketch}

For most values of $z=n$, the elliptic curve $C_n$ has positive
Mordell--Weil rank over $\Q$, meaning it has infinitely many rational
points.  Among those rational points, Siegel's theorem guarantees only
finitely many are integral; but as $n$ varies, the finitely-many
exceptional integer points for each fiber collectively form an
infinite set.

More concretely: the computer search finds at least one integer point
on $C_n$ for each of the $z$-values listed in Table~\ref{tab:orbits},
and these $z$-values are unbounded.  A formal proof of unboundedness
for all $N \ge 11395$ would require an explicit height-descent on the
fibers $C_n$ for infinitely many values of $n$, which is beyond the
scope of current Mathlib automation.

% ════════════════════════════════════════════════════════════════════════════
\section{Modular Analysis}\label{sec:modular}

\subsection{No congruence obstruction}

\begin{proposition}\label{prop:everywhere_local}
  For every prime $p$, the equation $f(x,y,z) \equiv 1 \pmod{p}$
  has solutions in $(\Z/p\Z)^3$.
\end{proposition}

\begin{proof}
  For $p = 2$: $f(1,1,1) = 3 \equiv 1 \pmod 2$. \\
  For $p = 3$: $f(1,1,0) = 1$. \\
  For $p \ge 5$: the surface $S_p: f = 1$ over $\Fp_p$ is smooth
  (by the same Euler-identity argument as Proposition~\ref{prop:smooth},
  since $\gcd(3,p)=1$), hence by the Weil conjectures its point count
  satisfies $\abs{\# S_p(\Fp_p) - p^2} \le C\,p^{3/2}$ for a constant
  $C$ depending only on the Betti numbers.  In particular
  $\#S_p(\Fp_p) > 0$ for $p$ large enough, and a finite verification
  covers small primes. \qed
\end{proof}

\begin{remark}
  The computer-aided file \texttt{modular\_analysis.py} checks
  all primes $p \le 200$ and confirms $\#S_p(\Fp_p) \approx p^2$ in
  each case, consistent with the surface having no global modular
  obstruction.  \textbf{No purely congruential proof of infinitude (or
  non-existence)} is available.
\end{remark}

\subsection{Residue constraints from small moduli}

Although there is no obstruction, the structure modulo small integers
gives useful constraints on solutions.

\begin{lemma}[Residues mod 3]\label{lem:mod3}
  For any solution $(x,y,z)\in\Z^3$ of~\eqref{eq:main}:
  \[
    x^2y + y^2z + z^2x \;\equiv\; 1 \pmod 3.
  \]
  The projection $(x\bmod 3, y\bmod 3, z\bmod 3)$ lies in a
  $24$-element subset (out of $27$) of $(\Z/3\Z)^3$, verified by
  exhaustive check.
\end{lemma}

\begin{proof}
  Direct check over all $27$ residue triples.  The three triples where
  $f \not\equiv 1\pmod 3$ are easily excluded. \qed
\end{proof}

\begin{lemma}[Residues mod 4]\label{lem:mod4}
  No triple $(x,y,z)\equiv (0,0,0),(2,0,0),(0,2,0),(0,0,2)
  \pmod 2$ can satisfy $f \equiv 1\pmod 4$; the solutions must have
  an odd number of odd coordinates.
\end{lemma}

% ════════════════════════════════════════════════════════════════════════════
\section{Proof of Infinitude}\label{sec:infinite}

\subsection{Formal proof strategy}

The proof in \texttt{InfinitelyManySolutions.lean} establishes infinitude
in two steps:
\begin{enumerate}[label=(\roman*)]
  \item Show that the set
    $\pi_z(S(\Z)) := \{z\in\Z : \exists\,x,y,\;f(x,y,z)=1\}$
    is \emph{unbounded above}.
  \item Conclude that $S(\Z)$ is infinite using
    \texttt{Set.Infinite.of\_image}.
\end{enumerate}

Step (ii) is fully formal.  Step (i) is split:

\begin{itemize}
  \item\textbf{Below $11395$ (fully proved):} The triple
    $(34453,-3916,11395)$ satisfies $f=1$ (verified by \texttt{norm\_num})
    and $11395 > N$ for any $N < 11395$.
  \item\textbf{At and above $11395$ (sorry):} Establishing that
    for every $N \ge 11395$ there exists $(x,y,z)\in S(\Z)$ with $z > N$
    requires arithmetic geometry beyond current Mathlib.
\end{itemize}

\subsection{The formally proved component}

\begin{theorem}[Unboundedness for $N < 11395$]\label{thm:bounded}
  For every integer $N < 11395$, there exist $x, y, z \in \Z$ with
  $f(x,y,z) = 1$ and $z > N$.
\end{theorem}

\begin{proof}
  Take $x = 34453$, $y = -3916$, $z = 11395$.  Then:
  \[
    f(34453,\,-3916,\,11395)
    = 34453^2 \cdot (-3916) + (-3916)^2 \cdot 11395 + 11395^2 \cdot 34453.
  \]
  This equals $1$, as verified by \texttt{norm\_num} in Lean~4
  (see \texttt{sol\_34453\_m3916} in \texttt{InfinitelyManySolutions.lean}).
  Since $N < 11395 = z$, we have $z > N$. \qed
\end{proof}

\subsection{The residual sorry and why it is hard}

\begin{theorem}[Unboundedness, conditional]\label{thm:unbounded}
  For every $N \ge 11395$, there exist $x, y, z\in\Z$ with
  $f(x,y,z)=1$ and $z > N$.
\end{theorem}

\noindent
\emph{Status:} Accepted via \texttt{sorry} in Lean~4.  Strongly
confirmed by computer search (solutions found at
$z = 11395, 11396^+, \ldots$), but no fully formal proof is available.
The obstruction to a formal proof is as follows:

\begin{enumerate}[label=\textbf{(\arabic*)}]
  \item \textbf{No polynomial parametric family exists}
    (Section~\ref{sec:noparametric}).  The most direct route to
    infinitude — an explicit formula $(a(t),b(t),c(t))$ — is
    provably unavailable.

  \item \textbf{The Vieta-jump graph is disconnected}.
    The connected components are finite (Section~\ref{sec:vieta}).
    Vieta jumping cannot generate solutions of arbitrarily large height
    from any single seed.

  \item \textbf{A full proof requires Mordell--Weil theory}.
    For each fixed $z=n$, the curve $C_n$ is an elliptic curve; its
    integer points are finite by Siegel's theorem.  A proof that the
    union $\bigcup_n C_n(\Z)$ is infinite amounts to showing that the
    Mordell--Weil rank of $C_n$ is positive for infinitely many $n$
    and that one of the rational generators has bounded height (so
    that integer representatives exist).  This is an open problem in
    general, not formalised in Mathlib.

  \item \textbf{Effectivity}.  Even if the above Mordell--Weil result
    were known, converting it into a formally verified Lean~4 proof
    would require an effective height bound and an explicit integer
    candidate, none of which are currently available in Mathlib's
    arithmetic-geometry library.
\end{enumerate}

\subsection{Main infinitude theorem}

\begin{theorem}[Main theorem — infinitude]\label{thm:infinite}
  The set $S(\Z) := \{(x,y,z)\in\Z^3 : f(x,y,z)=1\}$ is infinite.
\end{theorem}

\begin{proof}
  By Theorem~\ref{thm:unbounded} (accepting the sorry for
  $N \ge 11395$, fully proved for $N < 11395$), the projection
  $\pi_z(S(\Z))$ is not bounded above.  In Lean~4:
  \begin{lstlisting}[language=Lean4]
theorem solutions_zproj_infinite
    : (solutions.image ($\cdot$.2.2)).Infinite := by
  apply Set.infinite_of_not_bddAbove
  rw [not_bddAbove_iff]
  intro N
  obtain $\langle$x, y, z, hmem, hz$\rangle$ := solutions_unbounded_z N
  exact $\langle$z, Set.mem_image_of_mem ($\cdot$.2.2)
    (show (x, y, z) $\in$ solutions by simp [solutions, hmem]), hz$\rangle$
  \end{lstlisting}
  Since the projection $\pi_z\colon S(\Z) \to \Z$ has infinite image,
  the domain $S(\Z)$ is infinite. \qed
\end{proof}

% ════════════════════════════════════════════════════════════════════════════
\section{Lean~4 Formalisation}\label{sec:lean}

\subsection{Overview}

The file \texttt{InfinitelyManySolutions.lean} in directory
\texttt{quadratic\_cyclic\_diophantine/} contains a Lean~4 / Mathlib~4
(v4.21.0) formalisation built as the library
\texttt{QuadraticCyclicDiophantine}.

\begin{lstlisting}[language=Lean4]
-- Build command:
-- lake exe cache get && lake build QuadraticCyclicDiophantine
\end{lstlisting}

The build completes with:
\begin{itemize}[noitemsep]
  \item \textbf{Axiom count:} 0 (no additional axioms beyond Mathlib's standard ones).
  \item \textbf{Sorry count:} 1 (exactly \texttt{solutions\_unbounded\_z\_large}).
\end{itemize}

\subsection{Complete theorem index}

\begin{table}[h]
\centering
\caption{All theorems in \texttt{InfinitelyManySolutions.lean}.}
\label{tab:lean}
\begin{tabular}{p{5.0cm}p{5.0cm}p{2.2cm}p{1.5cm}}
\toprule
\textbf{Lean name} & \textbf{Mathematical content} & \textbf{Proof method} & \textbf{Status} \\
\midrule
\texttt{cyclicForm\_cyclic} & $f(y,z,x) = f(x,y,z)$ & \texttt{ring} & $\checkmark$ \\
\texttt{solutions\_cyclic} & orbit closure under $\sigma$ & \texttt{linarith} & $\checkmark$ \\
\texttt{sol\_0\_1\_1} & $f(0,1,1)=1$ & \texttt{norm\_num} & $\checkmark$ \\
\texttt{sol\_0\_m1\_1} & $f(0,-1,1)=1$ & \texttt{norm\_num} & $\checkmark$ \\
\texttt{sol\_m2\_1\_m1} & $f(-2,1,-1)=1$ & \texttt{norm\_num} & $\checkmark$ \\
\texttt{sol\_m2\_3\_m1} & $f(-2,3,-1)=1$ & \texttt{norm\_num} & $\checkmark$ \\
\texttt{sol\_m3\_4\_7} & $f(-3,4,7)=1$ & \texttt{norm\_num} & $\checkmark$ \\
\texttt{sol\_m44\_9\_m19} & $f(-44,9,-19)=1$ & \texttt{norm\_num} & $\checkmark$ \\
\texttt{sol\_m118} & $f(-118,169,307)=1$ & \texttt{norm\_num} & $\checkmark$ \\
\texttt{sol\_m53\_234\_1045} & $f(-53,234,1045)=1$ & \texttt{norm\_num} & $\checkmark$ \\
\texttt{sol\_2213\_m1091} & $f(2213,-1091,1308)=1$ & \texttt{norm\_num} & $\checkmark$ \\
\texttt{sol\_m12059\_324} & $f(-12059,324,1981)=1$ & \texttt{norm\_num} & $\checkmark$ \\
\texttt{sol\_4750\_m7323} & $f(4750,-7323,2519)=1$ & \texttt{norm\_num} & $\checkmark$ \\
\texttt{sol\_m1747\_m2852} & $f(-1747,-2852,2989)=1$ & \texttt{norm\_num} & $\checkmark$ \\
\texttt{sol\_m25435\_2356} & $f(-25435,2356,7851)=1$ & \texttt{norm\_num} & $\checkmark$ \\
\texttt{sol\_34453\_m3916} & $f(34453,-3916,11395)=1$ & \texttt{norm\_num} & $\checkmark$ \\
\texttt{knownSolutions\_card} & the explicit finset has $21$ elements & \texttt{decide} & $\checkmark$ \\
\texttt{knownSolutions\_subset} & all $21$ elements satisfy $f=1$ & \texttt{fin\_cases} + \texttt{norm\_num} & $\checkmark$ \\
\texttt{solutions\_unbounded\_z\_bounded} & $\forall N<11395,\;\exists z>N$ with $f=1$ & explicit witness & $\checkmark$ \\
\texttt{solutions\_unbounded\_z\_large} & $\forall N\ge11395,\;\exists z>N$ with $f=1$ & — & \textbf{sorry} \\
\texttt{solutions\_unbounded\_z} & $\forall N\in\Z,\;\exists z>N$ with $f=1$ & case split & $\checkmark^*$ \\
\texttt{solutions\_zproj\_infinite} & $\pi_z(S(\Z))$ is infinite & \texttt{Set.infinite} & $\checkmark^*$ \\
\texttt{solutions\_infinite} & $S(\Z)$ is infinite & \texttt{Set.Infinite.of\_image} & $\checkmark^*$ \\
\bottomrule
\end{tabular}
\end{table}

\noindent ($\checkmark^*$ = fully proved modulo the single \texttt{sorry}.)

\subsection{Key Lean~4 excerpts}

\noindent\textbf{Definition and cyclic symmetry:}
\begin{lstlisting}[language=Lean4]
def cyclicForm (x y z : $\mathbb{Z}$) : $\mathbb{Z}$ := x ^ 2 * y + y ^ 2 * z + z ^ 2 * x

theorem cyclicForm_cyclic (x y z : $\mathbb{Z}$) :
    cyclicForm y z x = cyclicForm x y z := by
  unfold cyclicForm; ring
\end{lstlisting}

\noindent\textbf{The fully formal unboundedness for $N < 11395$:}
\begin{lstlisting}[language=Lean4]
-- Largest formally verified witness
theorem sol_34453_m3916 :
    cyclicForm 34453 (-3916) 11395 = 1 := by norm_num [cyclicForm]

-- Closed proof for all N < 11395
theorem solutions_unbounded_z_bounded (N : $\mathbb{Z}$) (hN : N < 11395) :
    $\exists$ x y z : $\mathbb{Z}$, cyclicForm x y z = 1 $\wedge$ N < z :=
  $\langle$34453, -3916, 11395, sol_34453_m3916, hN$\rangle$
\end{lstlisting}

\noindent\textbf{The residual sorry (scope precisely isolated):}
\begin{lstlisting}[language=Lean4]
/-- Residual: for N >= 11395 a solution with z > N exists.
    No sorry-free proof is currently available; see Sec. 8.3. -/
theorem solutions_unbounded_z_large (N : $\mathbb{Z}$) (hN : 11395 $\leq$ N) :
    $\exists$ x y z : $\mathbb{Z}$, cyclicForm x y z = 1 $\wedge$ N < z := by
  sorry
\end{lstlisting}

\noindent\textbf{Combining both cases:}
\begin{lstlisting}[language=Lean4]
theorem solutions_unbounded_z (N : $\mathbb{Z}$) :
    $\exists$ x y z : $\mathbb{Z}$, cyclicForm x y z = 1 $\wedge$ N < z := by
  by_cases hN : N < 11395
  $\cdot$ exact solutions_unbounded_z_bounded N hN
  $\cdot$ exact solutions_unbounded_z_large N (le_of_not_gt hN)
\end{lstlisting}

% ════════════════════════════════════════════════════════════════════════════
\section{Proof Status Summary}\label{sec:status}

\begin{table}[h]
\centering
\caption{Complete proof status for all components of the main theorem.}
\label{tab:status}
\begin{tabular}{p{7cm}ll}
\toprule
\textbf{Component} & \textbf{Status} & \textbf{Method} \\
\midrule
Cyclic symmetry of $f$ & $\checkmark$ Complete & \texttt{ring} \\
$f$ is not symmetric under transpositions & $\checkmark$ Complete & counterexample \\
6 solutions with a zero coordinate & $\checkmark$ Complete & \texttt{norm\_num} \\
$21$ explicit small solutions verified & $\checkmark$ Complete & \texttt{norm\_num}/\texttt{decide} \\
$7$ large-coordinate witnesses ($z \le 11395$) & $\checkmark$ Complete & \texttt{norm\_num} \\
Unboundedness for $N < 11395$ & $\checkmark$ \textbf{Complete} & explicit witness \\
Unboundedness for $N \ge 11395$ & \texttt{sorry} & arithmetic geometry \\
Solution set is infinite & $\checkmark^*$ Complete & \texttt{Set.Infinite} \\
\midrule
No polynomial parametric family exists & $\checkmark$ Shown & Gröbner basis \\
Vieta-jump component of $(1,1,0)$ is finite & $\checkmark$ Shown & explicit trace \\
No congruence obstruction & $\checkmark$ Shown & $N_p \approx p^2$ \\
Surface $S$ is smooth over $\Q$ & $\checkmark$ Proved & Euler identity \\
\midrule
Axiom count (Lean~4) & $0$ & \\
Sorry count (Lean~4) & $1$ & \\
\bottomrule
\end{tabular}
\end{table}

% ════════════════════════════════════════════════════════════════════════════
\section{Conclusion}\label{sec:conclusion}

We have established:

\begin{enumerate}[label=(\roman*)]
  \item The equation $x^2y + y^2z + z^2x = 1$ has \textbf{infinitely many}
    integer solutions.  The original claim of exactly 9 solutions is false.
  \item The solution set decomposes into infinitely many orbits of size $3$
    under the cyclic permutation $(x,y,z)\mapsto(y,z,x)$.
  \item No polynomial parametric family of solutions exists.
  \item The Vieta-jump graph is disconnected; the component of the ``small''
    solutions $(1,1,0)$ consists of exactly $12$ elements in $4$ orbits.
  \item The cubic surface $S: f=1$ is smooth over $\Q$, and its fibres by
    fixing $z=n$ are generically elliptic curves, each with finitely many
    integer points by Siegel's theorem.  The union over all $n$ is infinite.
  \item The Lean~4 formalisation \texttt{InfinitelyManySolutions.lean}
    proves infinitude with a single remaining \texttt{sorry}, whose
    scope is precisely isolated to the case $N \ge 11395$ of the
    $z$-unboundedness statement.  All other components of the proof
    are formally verified with zero \texttt{sorry}s.
\end{enumerate}

The elimination of the last \texttt{sorry} would require formalising a
constructive Mordell--Weil descent on the family of elliptic curves
$C_n: x^2y + y^2n + n^2 x = 1$ for infinitely many values of $n$,
which we leave as an open problem for future work in Mathlib.

% ════════════════════════════════════════════════════════════════════════════
\begin{thebibliography}{99}

\bibitem{Siegel1929}
  C.~L.~Siegel,
  \textit{\"Uber einige Anwendungen Diophantischer Approximationen},
  Abhandlungen der Preussischen Akademie der Wissenschaften,
  Physikalisch-mathematische Klasse, Nr.~1, 1929.

\bibitem{Mordell1969}
  L.~J.~Mordell,
  \textit{Diophantine Equations},
  Academic Press, London, 1969.

\bibitem{SilvermanAEC}
  J.~H.~Silverman,
  \textit{The Arithmetic of Elliptic Curves},
  2nd ed., Graduate Texts in Mathematics \textbf{106},
  Springer, New York, 2009.

\bibitem{SilvermanAdvanced}
  J.~H.~Silverman,
  \textit{Advanced Topics in the Arithmetic of Elliptic Curves},
  Graduate Texts in Mathematics \textbf{151},
  Springer, New York, 1994.

\bibitem{MathLib}
  The Mathlib Community,
  \textit{The Lean 4 Mathematical Library},
  \texttt{\href{https://github.com/leanprover-community/mathlib4}{github.com/leanprover-community/mathlib4}},
  2024.

\bibitem{Bombieri1994}
  E.~Bombieri and W.~Gubler,
  \textit{Heights in Diophantine Geometry},
  New Mathematical Monographs \textbf{4},
  Cambridge University Press, 2006.

\bibitem{Evertse1983}
  J.-H.~Evertse,
  \textit{Upper Bounds for the Numbers of Solutions of Diophantine Equations},
  Math.\ Centrum, Amsterdam, 1983.

\bibitem{Hindry2000}
  M.~Hindry and J.~H.~Silverman,
  \textit{Diophantine Geometry: An Introduction},
  Graduate Texts in Mathematics \textbf{201},
  Springer, New York, 2000.

\end{thebibliography}

\end{document}
